\section*{Bab \ref{ch:g10:drainage-basins} --- Daerah Aliran Sungai, Sungai, dan Banjir}

\begin{solution}{exo:g10:drainage-basins:1}
\omterm{def:g10:drainage-basins:drainage-basin}{Daerah aliran sungai} untuk suatu titik di sungai adalah seluruh wilayah
daratan yang air lariannya mengalir ke sungai di hulu titik itu dan
melewatinya. \omterm{def:g10:drainage-basins:drainage-basin}{Garis pemisah air} adalah garis, di sepanjang punggung-punggung,
yang memisahkan dua daerah aliran yang bertetangga. Dua tetes yang jatuh
berjarak satu meter di kedua sisi \omterm{def:g10:drainage-basins:drainage-basin}{garis pemisah air} mengalir menuruni lereng
yang berlawanan, ke dalam dua daerah aliran yang berbeda: sungai-sungainya
dapat mencapai laut yang berbeda.
\end{solution}

\begin{solution}{exo:g10:drainage-basins:2}
\omterm{def:g10:drainage-basins:discharge}{Debit} adalah volume air yang melewati sebuah penampang sungai setiap detik,
dalam \unit{m^3/s}. \qty{1}{m^3/s} adalah \qty{1000}{L} setiap detik; dalam
sehari, $1000 \times 86\,400 = 86.4$ juta liter.
\end{solution}

\begin{solution}{exo:g10:drainage-basins:3}
Sebelum banjir, sungai itu membawa sekitar \qty{330}{m^3/s} (di antara
sekitar \num{240} dan \qty{390}{m^3/s}). Puncaknya \qty{9231}{m^3/s} pada
21 Januari. Kenaikan itu bermula pada 18 Januari (\qty{362}{m^3/s}): kenaikan
itu berlangsung tiga hari.
\end{solution}

\begin{solution}{exo:g10:drainage-basins:4}
$1/50 = 0.02$, \qty{2}{\percent}. Tidak: tahun-tahun itu saling bebas;
peluangnya tetap \qty{2}{\percent} tahun ini, apa pun yang terjadi tahun lalu.
\end{solution}

\begin{solution}{exo:g10:drainage-basins:5}
Satu kaki kubik adalah $0.3048^3 \approx \qty{0.02832}{m^3}$, jadi
$593\,000 \times 0.02832 \approx \qty{16800}{m^3/s}$. Amazon membawa
$206\,000/16\,800 \approx 12$ kali lebih banyak.
\end{solution}

\begin{solution}{exo:g10:drainage-basins:6}
Luas pita: $1.5 \times 0.3 = 0.45$, $0.90$, $1.05$, dan \qty{0.60}{m^2}:
$A = \qty{3.0}{m^2}$. \omterm{def:g10:drainage-basins:discharge}{Debit}: $0.2 \times 0.45 = 0.09$, $0.5 \times 0.90
= 0.45$, $0.6 \times 1.05 = 0.63$, $0.3 \times 0.60 = 0.18$: $Q =
\qty{1.35}{m^3/s}$. Kecepatan rata-rata $1.35/3.0 = \qty{0.45}{m/s}$.
\end{solution}

\begin{solution}{exo:g10:drainage-basins:7}
$66 \times 10^6 \times 0.020 = \qty{1.32e6}{m^3}$. Dua hari sama dengan
$2 \times 86\,400 = \qty{172800}{s}$: $1.32 \times 10^6/172\,800 \approx
\qty{7.6}{m^3/s}$ rata-rata.
\end{solution}

\begin{solution}{exo:g10:drainage-basins:8}
Diurutkan: 120, 88, 65, 59, 51, 47, 42, 38, 35. Dengan $n = 9$, $T = 10/m$:
yang terbesar (120) memiliki $T = 10$ tahun, yang kedua (88) 5 tahun, yang
ketiga (65) sekitar 3.3 tahun. Median, 51, berperingkat 5 dan $T = 2$ tahun:
puncak itu disamai atau dilampaui satu tahun dari dua.
\end{solution}

\begin{solution}{exo:g10:drainage-basins:9}
\omterm{def:g10:drainage-basins:hydrograph}{Hidrograf} daerah aliran berkota naik lebih awal dan jauh lebih tinggi, lalu
turun dengan cepat: di lereng curam, atap, dan jalan, hanya sedikit air yang
meresap dan air larian mencapai anak sungai dengan cepat dan sekaligus; \omterm{def:g10:drainage-basins:hydrograph}{waktu
tundanya} singkat. \omterm{def:g10:drainage-basins:hydrograph}{Hidrograf} daerah aliran berhutan naik lebih lambat, sampai
puncak yang lebih rendah dan lebih lebar: tanah dan hutan menyerap sebagian
besar hujan dan melepaskannya perlahan, dan air mengalir perlahan di lereng
yang landai. \omterm{def:g10:drainage-basins:hydrograph}{Aliran dasarnya} tetap lebih tinggi setelah hujan badai.
\end{solution}

\begin{solution}{exo:g10:drainage-basins:10}
$9231 \times 86\,400 \approx \qty{7.98e8}{m^3}$, disebarkan ke
$29\,900 \times 10^6 = \qty{2.99e10}{m^2}$: sekitar \qty{0.027}{m},
\qty{27}{mm}. \omterm{def:g10:drainage-basins:hydrograph}{Aliran dasar} memberikan $330 \times 86\,400/2.99 \times 10^{10}
\approx \qty{0.001}{m}$, sekitar \qty{1}{mm} sehari: hari puncak membawa
sekitar 28 kali lebih banyak.
\end{solution}

\begin{solution}{exo:g10:drainage-basins:11}
$P = 1 - 0.9^{10} \approx 1 - 0.349 = 0.651$: sekitar \qty{65}{\percent}.
Banjir itu sama sekali tidak datang dengan peluang $0.9^{10} \approx 0.35$,
sekitar \qty{35}{\percent}.
\end{solution}

\begin{solution}{exo:g10:drainage-basins:12}
\omterm{def:g10:drainage-basins:hydrograph}{Hidrograf} itu berbentuk segitiga sepanjang tiga hari dan setinggi
\qty{1200}{m^3/s}: luasnya $\tfrac12 \times 1200 \times 3 \times 86\,400 \approx
\qty{1.56e8}{m^3}$, \num{156} juta meter kubik. Bagian di atas
\qty{800}{m^3/s} adalah segitiga yang lebih kecil dengan bentuk yang sama, tiga
kali lebih kecil dalam tinggi ($400$, bukan $1200$) dan dalam lama waktu (satu
hari, bukan tiga): $\tfrac12 \times 400 \times 86\,400 \approx \qty{1.73e7}{m^3}$,
\num{17} juta meter kubik. Itu kurang dari \num{60} juta yang dapat disimpan
waduk: ya.
\end{solution}

\begin{solution}{exo:g10:drainage-basins:13}
Bila penuh sampai ke tepiannya, sungai itu sedalam
$800/(100 \times 2) = \qty{4}{m}$. (1) Dataran banjir membawa
\qty{400}{m^3/s} selebar \qty{900}{m} dengan kecepatan \qty{0.5}{m/s}:
dalamnya $400/(900 \times 0.5) \approx \qty{0.9}{m}$. (2) Di antara tanggul:
sedalam $1200/(100 \times 2) = \qty{6}{m}$, \qty{2}{m} di atas tepiannya, lebih
dari dua kali kedalaman di dataran banjir tanpa tanggul.
\end{solution}

\begin{solution}{exo:g10:drainage-basins:14}
Kita menginginkan $(1 - 1/T)^{30} = 0.5$, jadi $1 - 1/T = 0.5^{1/30} \approx 0.97716$,
$1/T \approx 0.02284$ dan $T \approx 44$ tahun.
\end{solution}

\begin{solution}{exo:g10:drainage-basins:15}
Setiap persegi mengirim airnya ke tetangganya yang terendah (baris dari utara
ke selatan, kolom dari barat ke timur):
\begin{center}
\begin{tabular}{|c|c|c|c|}
\hline
keluaran & barat & barat & barat \\ \hline
utara & barat & barat & selatan \\ \hline
utara & utara & selatan & selatan \\ \hline
utara & timur & timur & keluaran \\ \hline
\end{tabular}
\end{center}
Dengan mengikuti panah-panahnya, persegi 47, 39, 44, 37, 46, dan 30 mengalir
ke titik keluaran tenggara: 6 persegi, \qty{0.06}{km^2}. Sepuluh persegi
lainnya mengalir ke titik keluaran barat laut: \qty{0.10}{km^2}. \omterm{def:g10:drainage-basins:drainage-basin}{Garis pemisah
air} berjalan dari tepi timur di antara 50 dan 47, lalu di antara 48 dan 47, di
antara 48 dan 44, di antara 49 dan 44, di antara 49 dan 46, dan di antara 51
dan 46 sampai ke tepi selatan.
\end{solution}

\begin{solution}{pb:g10:drainage-basins:1}
\textbf{1.} $11\,560 \times 1.609^2 \approx \qty{29900}{km^2}$.

\textbf{2.} \omterm{def:g10:drainage-basins:drainage-basin}{Garis pemisah air} daerah aliran di hulu stasiun itu.

\textbf{3.} $484\,000 \times 0.3048^3 \approx \qty{13700}{m^3/s}$.

\textbf{4.} $13\,705/3341 \approx 4.1$ kali.

\textbf{5.} $A = 400 \times 8 = \qty{3200}{m^2}$ dan $v = Q/A = 13\,705/3200
\approx \qty{4.3}{m/s}$.

\textbf{6.} Sekitar \qty{330}{m^3/s}.

\textbf{7.} \qty{9231}{m^3/s}, pada 21 Januari.

\textbf{8.} Tiga hari.

\textbf{9.} $9231 \times 86\,400/(29\,900 \times 10^6) \approx
\qty{0.027}{m}$: sekitar \qty{27}{mm}.

\textbf{10.} Air baru mencapai sungai: hujan lain yang jatuh di daerah aliran
(atau salju yang mencair) mengirim gelombang air larian kedua yang lebih
kecil.

\textbf{11.} Tujuh belas hari setelah puncak. Kenaikan berasal dari air larian
lereng-lereng di dekat anak sungai, yang tiba dengan cepat dan sekaligus; pada
sisi turun sungai masih menerima air dari bagian-bagian terjauh daerah aliran
dan air yang terkuras perlahan dari tanah.

\textbf{12.} $T = n + 1 = 96$ tahun.

\textbf{13.} $96/2 = 48$ tahun.

\textbf{14.} $96/3 = 32$ tahun.

\textbf{15.} $96/4 = 24$ tahun dan $96/5 = 19.2$ tahun. Tidak banyak: dengan
cara mana pun, \omterm{def:g10:drainage-basins:discharge}{debit} ini datang kembali kira-kira sekali dalam dua puluh tahun.

\textbf{16.} Lima tahun (1936, 1937, 1943, 1972, dan 1996): $5/95 \approx
0.053$, sekitar \qty{5}{\percent} setahun, dengan \omterm{def:g10:drainage-basins:return-period}{periode ulang} sekitar
$96/5 \approx 19$ tahun.

\textbf{17.} $1/100 = 0.01$, \qty{1}{\percent}.

\textbf{18.} $1 - 0.01 = 0.99$.

\textbf{19.} $0.99^{30} \approx 0.740$.

\textbf{20.} $1 - 0.740 = 0.26$: \textbf{peluang sekurang-kurangnya satu banjir
100 tahunan dalam 30 tahun adalah sekitar \qty{26}{\percent}}, satu dari empat.
\end{solution}
